% =================================================================
%  CSE 815 -- Machine Learning
%  Module 1, Week 2, Part 1
%  Multiple Features, Vectorization, Gradient Descent for Multiple
%  Linear Regression
% =================================================================
\documentclass[aspectratio=169,10pt]{beamer}
\usepackage{cse815}

\title[CSE 815 -- Week 2, Part 1]{Regression with Multiple Input Variables}
\subtitle{Multiple Features, Vectorization, and Gradient Descent in $n$ Dimensions}
\author{Rokan Uddin Faruqui\\ \small Professor, Department of CSE}
\institute{\small University of Chittagong \\ \texttt{rufaruqui@cu.ac.bd}}
\date[Week 2 / Part 1]{Week 2 $\cdot$ Part 1}

\begin{document}

\begin{frame}[plain,noframenumbering]
  \titlepage
\end{frame}

% ---------------------------------------------------------------
\begin{frame}{Where we are, and where this goes}
  \begin{columns}[T,onlytextwidth]
    \begin{column}{0.5\textwidth}
      \begin{block}{Week 1 gave us}
        \[ \fwb(x)=wx+b \]
        \[ J(w,b)=\frac{1}{2m}\sum_{i=1}^{m}\bigl(\fwb(\xii{i})-\yii{i}\bigr)^2 \]
        \[ w := w-\alpha\tfrac{\partial J}{\partial w},\quad b := b-\alpha\tfrac{\partial J}{\partial b} \]
      \end{block}
      \small One feature: floor area. But nobody prices a flat on area alone ---
      bedrooms, age and location all matter.
    \end{column}
    \begin{column}{0.46\textwidth}
      \begin{keyidea}[Today]
        \small
        \begin{enumerate}
          \item Multiple features and the notation for them 
          \item The model as a dot product  
          \item Vectorization, and why it is fast 
          \item Gradient descent for $n$ features %\clo{25 min}
          \item The normal equation, in passing %\clo{10 min}
        \end{enumerate}
      \end{keyidea}
      \small Nothing conceptually new: the same three-part recipe, with $w$
      promoted from a number to a vector.
    \end{column}
  \end{columns}
\end{frame}

% ===============================================================
\section{Multiple Features}

\begin{frame}{One feature was never enough}
  \footnotesize
  \begin{center}
  \renewcommand{\arraystretch}{1.25}
  \begin{tabular}{@{}rrrrr@{}}
    \toprule
    \textbf{Size} & \textbf{Bedrooms} & \textbf{Floors} & \textbf{Age} & \textbf{Price}\\
    $x_1$ (k sq ft) & $x_2$ & $x_3$ & $x_4$ (yr) & $y$ (lakh BDT)\\
    \midrule
    2.10 & 5 & 1 &  45 & 46.0\\
    1.42 & 3 & 2 &  40 & 33.2\\
    1.53 & 3 & 2 &  30 & 31.5\\
    0.85 & 2 & 1 &  36 & 17.8\\
    1.94 & 4 & 2 &  15 & 42.0\\
    \bottomrule
  \end{tabular}
  \end{center}
  \begin{columns}[T,onlytextwidth]
    \begin{column}{0.48\textwidth}
      \begin{defnbox}[Notation]
        \small
        \begin{itemize}
          \item $n$ --- number of features (here $n=4$).
          \item $\xii{i}$ --- the $i$-th training example, now a \emph{vector}.
          \item $\xjj{i}{j}$ --- feature $j$ of example $i$.
        \end{itemize}
      \end{defnbox}
    \end{column}
    \begin{column}{0.48\textwidth}
      \small
      From the table:
      \[ \xii{2} = \begin{bmatrix} 1.42 & 3 & 2 & 40\end{bmatrix},\qquad \xjj{2}{3}=2. \]
      $\xii{i}$ is a row of the table; $\xjj{i}{j}$ is one cell.
    \end{column}
  \end{columns}
\end{frame}

\begin{frame}{The model with $n$ features}
  \small
  Previously $\fwb(x)=wx+b$. Now every feature carries its own weight:
  \[
    f_{\vw,b}(\vx) \;=\; w_1x_1 + w_2x_2 + \dots + w_nx_n + b
    \;=\; \vw\cdot\vx + b ,
  \]
  writing $\vw=[w_1,\dots,w_n]^{\!\top}$, $\vx=[x_1,\dots,x_n]^{\!\top}$ and
  $\vw\cdot\vx=\sum_{j=1}^{n}w_jx_j$ for the dot product. The bias $b$ stays a scalar.
  \begin{columns}[T,onlytextwidth]
    \begin{column}{0.52\textwidth}
      \footnotesize
      A plausible fitted model for our flats:
      \[ f(\vx) = 0.1x_1 + 4x_2 + 10x_3 - 2x_4 + 80 \]
      Every parameter reads as a price: $b=80$ is a base in lakh, $w_2=4$ says an extra
      bedroom adds 4 lakh, $w_4=-2$ says each year of age costs 2 lakh.
    \end{column}
    \begin{column}{0.44\textwidth}
      \begin{keyidea}
        \footnotesize
        This is \textbf{multiple} linear regression --- several \emph{inputs}. It is not
        ``multivariate regression'', which means several \emph{outputs}.
      \end{keyidea}
    \end{column}
  \end{columns}
\end{frame}

% ===============================================================
\section{Vectorization}

\begin{frame}{Three ways to compute the same thing}
  \footnotesize
  With $n=4$, $\vw = [w_1,\dots,w_4]$ and $\vx=[x_1,\dots,x_4]$:
  \begin{block}{1. Written out}
    \texttt{f = w[0]*x[0] + w[1]*x[1] + w[2]*x[2] + w[3]*x[3] + b}\\
    Fine for $n=4$; unwritable for $n=100$.
  \end{block}
  \begin{block}{2. A loop}
    \texttt{f = 0}\\
    \texttt{for j in range(n): f = f + w[j] * x[j]}\\
    \texttt{f = f + b}\\
    Scales in source length, but runs in the Python interpreter.
  \end{block}
  \begin{block}{3. Vectorized}
    \texttt{f = np.dot(w, x) + b}\\
    Shorter \emph{and} far faster --- the only version you should write.
  \end{block}
\end{frame}

\begin{frame}{Why vectorization is faster}
  \framesubtitle{It is not that the loop is ``interpreted'' --- it is that the hardware works in parallel}
  \begin{columns}[T,onlytextwidth]
    \begin{column}{0.48\textwidth}
      \centering {\scriptsize Sequential: one product per time step}\par
      \begin{tikzpicture}[font=\scriptsize,node distance=2mm]
        \foreach \i/\t in {0/$t_0$,1/$t_1$,2/$t_2$,3/$t_3$}{
          \node[draw,rounded corners,fill=culight,minimum width=1.5cm,minimum height=5mm]
            at (0,-\i*0.72) (n\i) {$w_\i x_\i$};
          \node[left=2mm of n\i,cugray] {\t};
        }
        \node[draw,fill=cuamber!25,rounded corners,minimum width=1.5cm,minimum height=5mm]
          at (0,-3.3) (s) {sum};
        \draw[-{Stealth},cugray] (n3) -- (s);
      \end{tikzpicture}
    \end{column}
    \begin{column}{0.48\textwidth}
      \centering {\scriptsize Vectorized: all products in one time step}\par
      \begin{tikzpicture}[font=\scriptsize]
        \foreach \i in {0,1,2,3}{
          \node[draw,rounded corners,fill=cuteal!20,minimum width=1.3cm,minimum height=5mm]
            at (\i*1.45,0) (v\i) {$w_\i x_\i$};
        }
        \node[cugray] at (-1.0,0) {$t_0$};
        \node[draw,fill=cuamber!25,rounded corners,minimum width=1.5cm,minimum height=5mm]
          at (2.2,-1.4) (s) {sum};
        \foreach \i in {0,1,2,3}{\draw[-{Stealth},cugray] (v\i) -- (s);}
        \node[cugray] at (-1.0,-1.4) {$t_1$};
      \end{tikzpicture}
    \end{column}
  \end{columns}
  \vspace{2mm}
  \footnotesize
  \texttt{np.dot} hands the arrays to compiled code using the CPU's \textbf{SIMD}
  instructions --- one instruction acting on several numbers at once --- then a parallel
  tree reduction for the sum. On a GPU the same idea spans thousands of cores.
  \begin{keyidea}
    The speed-up grows with $n$ and $m$, so vectorization is a correctness-level concern,
    not a style preference.
  \end{keyidea}
\end{frame}

\begin{frame}{Vectorizing the gradient descent step}
  \footnotesize
  Suppose $n=16$ and we have all sixteen derivatives $d_1,\dots,d_{16}$.
  \vspace{1mm}
  \begin{columns}[T,onlytextwidth]
    \begin{column}{0.48\textwidth}
      \begin{block}{Without vectorization}
        \texttt{for j in range(16):}\\
        \texttt{\ \ \ \ w[j] = w[j] - 0.1 * d[j]}
        \vspace{1mm}
        Sixteen steps, executed one after another.
      \end{block}
    \end{column}
    \begin{column}{0.48\textwidth}
      \begin{block}{With vectorization}
        \texttt{w = w - 0.1 * d}
        \vspace{1mm}
        One step. NumPy updates all sixteen weights in parallel.
      \end{block}
    \end{column}
  \end{columns}
  \begin{pitfall}
    \small
    \texttt{w * x} is \emph{elementwise} and returns a vector; \texttt{np.dot(w, x)} sums
    those products and returns the scalar the model needs. Confusing them yields an array of
    the wrong shape, often with no error until much later.
  \end{pitfall}
\end{frame}

% ===============================================================
\section{Gradient Descent for Multiple Features}

\begin{frame}{The cost, and the gradient}
  With $\vw\in\R^n$ and $b\in\R$, the cost is unchanged in form:
  \[
    J(\vw,b) \;=\; \frac{1}{2m}\sum_{i=1}^{m}\bigl(f_{\vw,b}(\xii{i})-\yii{i}\bigr)^{2},
    \qquad f_{\vw,b}(\xii{i}) = \vw\cdot\xii{i}+b .
  \]
  Differentiating with respect to one weight $w_j$ --- everything else is a constant:
  \[
    \frac{\partial J}{\partial w_j}
      = \frac{1}{m}\sum_{i=1}^{m}\bigl(f_{\vw,b}(\xii{i})-\yii{i}\bigr)\,\xjj{i}{j},
    \qquad
    \frac{\partial J}{\partial b}
      = \frac{1}{m}\sum_{i=1}^{m}\bigl(f_{\vw,b}(\xii{i})-\yii{i}\bigr).
  \]
  \begin{keyidea}
    Compare with Week 1: $\partial J/\partial w = \frac1m\sum(\fwb(\xii{i})-\yii{i})\xii{i}$.
    The only change is that $\xii{i}$ became $\xjj{i}{j}$ --- feature $j$ of example $i$.
    The residual is shared by every parameter; only the weighting differs.
  \end{keyidea}
\end{frame}

\begin{frame}{The algorithm}
  \begin{defnbox}[Gradient descent, $n$ features]
    Repeat until convergence, updating all $n+1$ parameters \emph{simultaneously}:
    \[
      \begin{aligned}
        w_j &:= w_j - \alpha\,\frac{1}{m}\sum_{i=1}^{m}\bigl(f_{\vw,b}(\xii{i})-\yii{i}\bigr)\xjj{i}{j}
        \qquad \text{for } j = 1,\dots,n\\[2pt]
        b &:= b - \alpha\,\frac{1}{m}\sum_{i=1}^{m}\bigl(f_{\vw,b}(\xii{i})-\yii{i}\bigr)
      \end{aligned}
    \]
  \end{defnbox}
  \footnotesize
  In vector form the whole update is two lines:
  \[
    \vw := \vw - \alpha\,\nabla_{\vw}J, \qquad b := b - \alpha\,\frac{\partial J}{\partial b},
    \qquad
    \nabla_{\vw}J = \frac{1}{m}X^{\!\top}(X\vw + b\mathbf{1} - \mathbf{y}) .
  \]
  \begin{pitfall}
    ``Simultaneously'' now means all $n+1$ parameters. Compute the entire gradient from the
    old $(\vw,b)$, then assign. In NumPy: build \texttt{dj\_dw} and \texttt{dj\_db} first,
    and only then write \texttt{w -= alpha*dj\_dw; b -= alpha*dj\_db}.
  \end{pitfall}
\end{frame}

\begin{frame}{In code}
  \footnotesize
  \begin{block}{Computing the gradient for all $n$ features}
    \texttt{def compute\_gradient(X, y, w, b):}\\
    \texttt{\ \ \ \ m = X.shape[0]}\\
    \texttt{\ \ \ \ err   = (X @ w + b) - y\ \ \ \ \ \ \ \ \ \ \# shape (m,)}\\
    \texttt{\ \ \ \ dj\_dw = (X.T @ err) / m\ \ \ \ \ \ \ \ \ \ \# shape (n,)}\\
    \texttt{\ \ \ \ dj\_db = np.sum(err) / m\ \ \ \ \ \ \ \ \ \ \# scalar}\\
    \texttt{\ \ \ \ return dj\_dw, dj\_db}
  \end{block}
  \begin{columns}[T,onlytextwidth]
    \begin{column}{0.48\textwidth}
      \textbf{Shapes are the contract}
      \begin{itemize}
        \item \texttt{X} is $(m,n)$, \texttt{y} is $(m,)$
        \item \texttt{w} is $(n,)$, \texttt{b} is a scalar
        \item \texttt{X @ w} is $(m,)$ --- one prediction per example
        \item \texttt{X.T @ err} is $(n,)$ --- one derivative per feature
      \end{itemize}
    \end{column}
    \begin{column}{0.48\textwidth}
      \textbf{Why \texttt{X.T @ err} is the gradient}\\[2pt]
      Row $j$ of $X^{\!\top}$ is feature $j$ across all examples. Dotting it with the residual
      vector gives $\sum_i(\hat y^{(i)}-y^{(i)})x^{(i)}_j$ --- exactly
      $m\,\partial J/\partial w_j$.
    \end{column}
  \end{columns}
\end{frame}

% ===============================================================
\section{An Alternative: the Normal Equation}

\begin{frame}{Solving for $\vw$ in one step}
  \footnotesize
  Linear regression is unusual: its minimum has a closed form. Setting $\nabla J=\mathbf{0}$
  and solving gives the \textbf{normal equation}
  \[
    \boxed{\;\vw = (X^{\!\top}X)^{-1}X^{\!\top}\mathbf{y}\;}
  \]
  where $X$ carries a column of ones to absorb $b$.
  \begin{columns}[T,onlytextwidth]
    \begin{column}{0.48\textwidth}
      \begin{block}{In its favour}
        \begin{itemize}
          \item No $\alpha$ to choose.
          \item No iterating, no convergence check.
          \item Exact, up to floating point.
        \end{itemize}
      \end{block}
    \end{column}
    \begin{column}{0.48\textwidth}
      \begin{alertblock}{Against}
        \begin{itemize}
          \item Only for \emph{linear} regression --- no logistic regression, no neural nets.
          \item Needs $(X^{\!\top}X)^{-1}$: cost grows like $n^3$, so it is unusable past
                $n\sim10^4$.
          \item Fails when $X^{\!\top}X$ is singular (duplicated or redundant features).
        \end{itemize}
      \end{alertblock}
    \end{column}
  \end{columns}
  \begin{keyidea}
    Use it as a \emph{test oracle}: on a small problem, check that your gradient descent
    converges to what \texttt{np.linalg.lstsq} returns. Gradient descent remains the method
    that generalises to everything else in this course.
  \end{keyidea}
\end{frame}

% ===============================================================
\section{Wrap-up and Assessment}

\begin{frame}{Summary of Part 1}
  \small
  \begin{columns}[T,onlytextwidth]
    \begin{column}{0.48\textwidth}
      \begin{block}{Concepts}
        \begin{itemize}
          \item $n$ features; $\xii{i}$ is a vector, $\xjj{i}{j}$ one of its entries.
          \item $f_{\vw,b}(\vx)=\vw\cdot\vx+b$.
          \item Vectorization is parallel hardware, not tidier code.
          \item The gradient's form is unchanged; only the weighting index moves.
          \item The normal equation solves this one model exactly, and nothing else.
        \end{itemize}
      \end{block}
    \end{column}
    \begin{column}{0.48\textwidth}
      \begin{block}{Formulas}
        \[ f_{\vw,b}(\vx)=\vw\cdot\vx+b \]
        \[ \frac{\partial J}{\partial w_j}=\frac1m\sum_{i=1}^{m}\bigl(f_{\vw,b}(\xii{i})-\yii{i}\bigr)\xjj{i}{j} \]
        \[ \vw := \vw - \alpha\nabla_{\vw}J \]
      \end{block}
      \begin{block}{Next}
        \small Part 2: why unscaled features cripple gradient descent, how to check
        convergence, choosing $\alpha$, and building new features.
      \end{block}
    \end{column}
  \end{columns}
\end{frame}

\begin{frame}{Descriptive questions}
  \footnotesize
  \begin{enumerate}
    \item Define $n$, $m$, $\xii{i}$ and $\xjj{i}{j}$. For the flats table on slide 4, state
          $\xii{4}$, $\xjj{4}{2}$ and $n$. \hfill\clo{5}
    \item Write the multiple linear regression model in both expanded and vector form, and
          explain why $b$ is not part of $\vw$. \hfill\clo{4}
    \item Distinguish \emph{multiple} linear regression from \emph{multivariate} regression.
          \hfill\clo{3}
    \item Derive $\partial J/\partial w_j$ for the multi-feature squared-error cost. State
          precisely where it differs from the single-feature case and why. \hfill\clo{8}
    \item Explain what vectorization means and why \texttt{np.dot(w,x)} outperforms an
          equivalent Python loop. Refer to SIMD in your answer. \hfill\clo{6}
  \end{enumerate}
\end{frame}

\begin{frame}{Descriptive questions (continued)}
  \footnotesize
  \begin{enumerate}\setcounter{enumi}{5}
    \item For $\vw=[0.1,4,10,-2]$ and $b=80$, compute the predicted price of a flat with
          $\vx=[1.5,3,2,20]$. Interpret each of the four terms in words. \hfill\clo{6}
    \item Given \texttt{X} of shape $(m,n)$, \texttt{w} of shape $(n,)$ and \texttt{y} of
          shape $(m,)$, state the shape of \texttt{X @ w}, \texttt{X.T @ err} and
          \texttt{np.sum(err)}. Explain why \texttt{X.T @ err} is the weight gradient.
          \hfill\clo{6}
    \item State the normal equation and give two reasons it does not replace gradient
          descent in this course. \hfill\clo{5}
    \item A student writes \texttt{f = w * x + b} instead of \texttt{np.dot(w,x) + b}.
          Describe what is actually computed, its shape, and why the bug may go unnoticed.
          \hfill\clo{5}
  \end{enumerate}
\end{frame}

\begin{frame}{Design questions}
  \footnotesize
  \begin{enumerate}
    \item \textbf{Design the feature set.} You must predict monthly electricity bills for
          households in Chattogram. Propose five features, giving units for each, and state
          the sign you expect each weight to take and why. Identify one feature that looks
          useful but would leak the answer. \hfill\clo{12}
    \item \textbf{Design for interpretability.} A bank will use your linear model to explain
          loan decisions to customers. Which choices in the model --- feature units, feature
          selection, the presence of $b$ --- affect whether $w_j$ can be quoted as a reason?
          Write the sentence you would put in front of a customer for one feature.
          \hfill\clo{10}
    \item \textbf{Design a shape-check harness.} Write three assertions that
          \texttt{compute\_gradient(X, y, w, b)} should satisfy for any valid input, such
          that a transposed matrix or a misused \texttt{*} is caught immediately. State what
          each assertion rules out. \hfill\clo{10}
    \item \textbf{Choose the solver.} For each case, argue for gradient descent or the normal
          equation: (a) $m=500$, $n=6$; (b) $m=10^7$, $n=3000$; (c) $n=50$ where two features
          are duplicates; (d) a logistic regression on $m=2000$, $n=10$. \hfill\clo{12}
  \end{enumerate}
\end{frame}

\begin{frame}{Reading}
  \begin{block}{Primary text}
    G.~James, D.~Witten, T.~Hastie and R.~Tibshirani,
    \emph{An Introduction to Statistical Learning}, 2nd ed., Springer, 2021.\\[2pt]
    \hfill{\small$\rightarrow$ \textbf{\S3.2} --- multiple linear regression: estimating the
    coefficients, interpreting them, and what the individual $w_j$ do and do not mean.}
  \end{block}
  \begin{block}{Supplementary}
    C.~M.~Bishop, \emph{Pattern Recognition and Machine Learning}, Springer, 2006.\\[2pt]
    \hfill{\small$\rightarrow$ \textbf{\S3.1--3.1.1} --- linear basis function models and the
    normal equations in matrix form.}
  \end{block}
  \vspace{2mm}
  \begin{center}
    \small\color{cugray}
    Lab (CSE 816, Week 2): NumPy vectorization; multiple linear regression.\\
    Practice quiz: \emph{Multiple linear regression}.
  \end{center}
\end{frame}

\end{document}
